\chapter{What Has Been Shown}
\label{ch:conclusion}

\section{The argument in brief}

The difficulty of watching certain films at home is usually described as a matter of loudness. This monograph has argued that the better description is occupation: the extent to which a soundtrack stays far above the level its dialogue requires, and for how long it does so without returning. Standard loudness measures summarise distributions and cannot see this temporal property. A listener anchored to dialogue meets it directly, as a sequence of exceedances to be managed by hand, and a long, deep block of exceedance leaves only one practical response, which is to turn the sound off and wait.

\section{Graded claims}

Table~\ref{tab:grades} collects the principal claims of the book with their grades.

\begin{longtable}{p{0.62\linewidth}lc}
\toprule
Claim & Grade & Chapter \\
\midrule
\endhead
Integrated loudness and loudness range are distributional and do not describe the temporal arrangement of exceedances & Established & \ref{ch:background} \\
Cinema reproduction is calibrated to a reference level; home reproduction generally is not & Established & \ref{ch:background} \\
Perceived control over noise reduces its aftereffects & Established & \ref{ch:background} \\
Speech intelligibility depends most on roughly the one to four kilohertz region & Established & \ref{ch:background} \\
The default \texttt{ffmpeg} mono downmix of a tagged 5.1 stream omits the LFE channel & Measured & \ref{ch:pipeline} \\
The film's relative recovery windows are ordinarily frequent (median gap 17.2 s) but absent for one stretch of 694.7 s & Measured & \ref{ch:case} \\
The film's level is high and sustained from roughly 61 to 80 minutes & Measured & \ref{ch:case} \\
Low-frequency onset activity recurs throughout the film rather than only at its climax & Measured & \ref{ch:case} \\
The film's integrated loudness is \resIntegrated~LUFS and its loudness range \resLRA~LU & Measured & \ref{ch:case} \\
The dialogue reference lies \xDialogueBelowIntegrated~dB below integrated loudness & Measured & \ref{ch:case} \\
\resExceedTen{} of audible time lies at least 10~dB above subtitled dialogue level, but only \resExceedTwenty{} at least 20~dB above; loud material forms a plateau from about $+12$ to $+19$~dB & Measured & \ref{ch:case} \\
The film contains \resOccBlocksN{} occupation blocks of at least a minute, six of them after the first hour & Measured & \ref{ch:case} \\
From \xSpanStart{} to \xSpanEnd{} the soundtrack returns to dialogue level for only \xSpanRecoverySeconds~s & Measured & \ref{ch:case} \\
Relative and absolute definitions of recovery locate the longest gap in the same place & Measured & \ref{ch:case} \\
The loud plateau reflects compression or limiting in the mix or master & Hypothesised & \ref{ch:case} \\
\resSubInBlocksShare{} of subtitled dialogue falls inside occupation blocks & Measured & \ref{ch:case} \\
The centre-channel proxy locates dialogue poorly (precision and recall near one half) but estimates its level within about 1~dB & Measured & \ref{ch:case} \\
Excluding exceedance itself, occupation blocks are not more overdetermined than the rest of the film & Measured & \ref{ch:overdetermination} \\
Most large swells begin from audible levels rather than near silence & Measured & \ref{ch:case} \\
Level-defined occupation blocks coincide with sustained transient barrages & Rejected & \ref{ch:case} \\
Volume reductions are predicted by occupation variables better than by level or whole-film statistics (H1) & Hypothesised & \ref{ch:framework} \\
Reduction probability rises with recovery debt (H2) & Hypothesised & \ref{ch:framework} \\
Reduction probability rises with overdetermination at fixed exceedance (H3) & Hypothesised (premise weakened) & \ref{ch:overdetermination} \\
The bass-amplifier chain widened the dialogue-to-action gap and added distortion to loud passages (H4) & Hypothesised & \ref{ch:chain} \\
A dialogue-forward downmix reduces throttling without loss of comprehension (H5) & Hypothesised & \ref{ch:intervention} \\
Version 3 swell candidates identify swells & Rejected & \ref{ch:failures} \\
Version 3 harshness peaks identify harsh material & Rejected & \ref{ch:failures} \\
The stretch from 81 to 84 minutes is the most intense in the film (it is the quiet closing dialogue scene) & Rejected & \ref{ch:failures} \\
Percentile-defined totals show that the film's recovery is scarce & Rejected & \ref{ch:failures} \\
The early analysis had established that the film ``has plenty of dynamic range'' & Rejected & \ref{ch:failures} \\
Swell events of 30 dB or more are musical crescendos & Rejected (untested) & \ref{ch:failures} \\
The low-frequency detector identifies drums & Rejected & \ref{ch:failures} \\
\bottomrule
\caption{Principal claims and their grades.}\label{tab:grades}
\end{longtable}

\section{Limitations}

The study rests on one film, one listener, one viewing, and one unmeasured reproduction chain. Its dialogue reference depends on subtitle timing, which marks when speech occurs but not every utterance, and which may be imperfectly aligned with the analysed release. Its frame-level detectors, apart from the swell detector, still use unweighted features and film-internal percentile thresholds, even though the occupation measures now use BS.1770 loudness. Its residualised overdetermination index removes only the smooth dependence of each channel on level. Its response data consist of a retrospective report without timestamps. Each of these limitations has a remedy specified in the text: pooled thresholds across a corpus, psychoacoustic replacements for the spectral descriptors, source labels from blinded annotation, the logging protocol of Chapter~\ref{ch:response}, and the chain measurements of Chapter~\ref{ch:chain}.

\section{A programme}

The research programme that follows from the framework has five components, in order of cost:
\begin{enumerate}
\item characterise the reproduction chain with a sine sweep and a level calibration;
\item conduct blind, multi-listener annotation of the audition clips with random controls;
\item log volume behaviour for new viewings under at least two playback conditions, including the dialogue-forward downmix;
\item analyse a comparison corpus with common, dialogue-referenced thresholds;
\item fit the hazard models of Chapter~\ref{ch:response} and report which hypotheses survive.
\end{enumerate}

\section{Closing}

The instruments built for this study failed several times before they measured anything, and the failures were more convincing than the successes until they were checked. That history is part of the result. A complaint about a soundtrack is easy to confirm with plots that are drawn after the complaint is made. It is harder, and more useful, to specify in advance what the soundtrack would have to contain for the complaint to be explained, to measure whether it does, and to name the listener data that could show the explanation wrong. The concept of occupation is offered as the core of such a specification.
